
\subsubsection{6.1 Key Findings}

\textbf{The LC correction preserves capability ordering.} r\_partial = 0.985 provides empirical validation of AlpacaEval 2.0's design goal across 223 diverse models spanning GPT-4 class to lightweight instruction-tuned variants, as operationalized by human preference win\_rate.

\textbf{Why partial correlation exceeds bivariate.} Verbosity acts as a mild suppressor ($\rho$(win\_rate, avg\_length) inferred from OLS $R^{2}(win_rate$ ~ avg\_length) = 0.433 in H-M2 residualization; Pearson R $\approx$ 0.658, with Spearman $\rho$ approximately 0.63). Partialling out verbosity reveals the capability signal more clearly, strengthening from the present sample bivariate of 0.966 to 0.985. This is consistent with Hu et al. [2024]: the desirability (length-independent) channel in win\_rate survives LC correction without attenuation.

\textbf{The verbosity penalty.} $\beta _len = -4.37$ (negative) confirms the LC correction penalizes verbosity once capability is controlled --- exactly the intended behavior of the debiasing procedure.

\subsubsection{6.2 The Alignment Gap Nuance}

The non-monotonic Q3 median $\Delta$ = 5.43 > Q4 = 0.91 reflects two factors: (1) mathematical dependency ($\Delta$ = LC\_winrate - win\_rate; grouping by win\_rate quartiles makes $\Delta$ partly self-referential, as confirmed by Spearman(win\_rate, $\Delta ) = -0.050$, p = 0.455); and (2) possible heterogeneity in verbosity-exploitation strategies within Q3, where mid-high capability models may produce verbose outputs that are most substantially adjusted by the LC correction. The $10\times$ effect-size difference between H-M3 ($\epsilon ^2$ = 0.088) and H-C1 ($\epsilon ^2$ = 0.883) establishes LC\_winrate as the more appropriate dependent variable for capability-alignment studies. Using $\Delta$ as dependent variable conflates the mathematical composition structure with the empirical signal and substantially attenuates observed effects.

\subsubsection{6.3 Limitations}

\textbf{L1 --- Observational study}: The analyses establish association, not causation. Unmeasured confounders (model family, RLHF budget, base model, training data composition) could co-vary with both metrics. Causal attribution requires experimental designs not feasible for N = 223 heterogeneous models.

\textbf{L2 --- Single dataset}: Results derive from AlpacaEval 2.0 (predominantly 2023--2024 era models). Generalization to MT-Bench, MMLU, Chatbot Arena, or post-2024 frontier models is untested. Cross-framework replication is identified as immediate future work.

\textbf{L3 --- Imperfect capability proxy}: win\_rate conflates capability with response style preferences and annotator biases. All capability-related interpretations should be understood as ''win\_rate independently predicts LC preference beyond verbosity'' rather than making strong claims about model cognitive capability. The operationalization is supported by Dubois et al.'s [2024] finding that win\_rate correlates with Chatbot Arena ELO ($\rho \approx$ 0.98).

\textbf{L4 --- Mild heteroscedasticity}: Breusch-Pagan p = 0.011 in H-M1. OLS $\beta$ estimates remain unbiased; standard errors may be slightly imprecise. Non-critical given p\_win = 4.58e-145.

\textbf{L5 --- $\Delta -monotonicity$ not confirmed}: Step 3 of the original hypothesis is only partially supported. The bidirectional alignment gap varies significantly across capability levels (medium effect, KW p = 5.97e-05) but does not follow a strictly monotonic pattern. LC\_winrate-based evidence (H-C1, large effect) provides stronger support for the capability-alignment relationship.

\subsubsection{6.4 Implications for Evaluation Study Design}

The $10\times$ effect-size attenuation from choosing $\Delta$ over LC\_winrate as dependent variable has direct implications for study design. Researchers analyzing capability-LC relationships should prefer LC\_winrate as the primary outcome rather than the alignment gap metric, which introduces mathematical dependency when models are grouped by win\_rate. This distinction is empirically demonstrated here rather than merely argued on theoretical grounds.

